\subsection{Deliverable 3: Delta Sheet Summarising Differences Between Stated and Observed Practice}
\label{sec:deliverable3}

The comparison maps interview activities (A01-A09) to observed events in Deliverable 2. A delta was registered when at least one of the following differed: event presence, order, timing, workaround behavior, or interpretation of activity labels. Methodologically, this follows Chapter 27's logic of comparing observed process traces with a pre-existing process description to identify frequent exceptions and model limits \parenciteshort{guidehealthinformatics2026ch27}.

Table~\ref{tab:d3_delta_sheet} summarises the full delta comparison between stated and observed routine. Key deltas are:
\begin{itemize}
  \item \textit{Missing event}: no distinct seated breakfast-and-schedule-review step in observation.
  \item \textit{Extra event}: additional roommate coordination event appears only in observed data.
  \item \textit{Different order}: bag packing occurs after breakfast preparation begins.
  \item \textit{Timing distortion}: hygiene block starts later than stated baseline.
  \item \textit{Workaround/coordination step}: missing lock triggers ad hoc borrowing via messaging.
  \item \textit{Relabeling/ambiguity}: ``breakfast'' shifts from sit-down meal to banana plus coffee-to-go.
\end{itemize}

All six required categories are documented in Table~\ref{tab:d3_delta_sheet}: missing events, extra events, different order, timing distortions, workarounds/coordination steps, and relabeling/ambiguity.

Overall, the \textit{interview-based DCR} (Deliverable 1) captures \textit{knowledge-as-stated} and \textit{work-as-imagined}, while the observed event sequence (\textit{observed DCR} evidence) captures \textit{knowledge-as-discovered}, \textit{work-as-done}, and \textit{work-as-recorded}. The delta sheet therefore provides traceable evidence of the \textit{say-do gap}.

\begin{table}[H]
\centering
\footnotesize
\caption{Delta sheet between interview model and observed routine}
\label{tab:d3_delta_sheet}
\rowcolors{3}{black!4}{white}
\begin{tabularx}{\textwidth}{
>{\raggedright\arraybackslash}p{0.08\textwidth}
>{\raggedright\arraybackslash}p{0.18\textwidth}
>{\raggedright\arraybackslash}p{0.15\textwidth}
>{\raggedright\arraybackslash}p{0.15\textwidth}
>{\raggedright\arraybackslash}X
>{\raggedright\arraybackslash}p{0.16\textwidth}
}
\toprule
Delta ID & Category & Interview ref. & Observed ref. & Difference description & Implication \\
\midrule
D01 & Missing events & A06 & No direct event & No distinct seated breakfast-and-schedule-review event occurred. & Interview model overstates routine regularity. \\
D02 & Extra events & None in baseline & E03 (06:52) & Short roommate-message coordination appears in observed routine only. & Reveals practical coordination load. \\
D03 & Different order & A04 before A05 & E07 then E09 & Bag packing happened after breakfast preparation started. & Ordering constraints are weaker than stated. \\
D04 & Timing distortions & A03 expected 06:50-07:02 & E05 start 06:57 & Hygiene block shifted later due to wake-up delays. & Time windows in interview are optimistic. \\
D05 & Workarounds/coordination steps & A08 final check & E10 (07:21) & Missing lock triggered ad hoc borrowing via message. & Recovery actions are underrepresented in interview model. \\
D06 & Relabeling/ambiguity & A05/A06 & E08, E13 & ``Breakfast'' changed from sit-down meal to banana plus coffee-to-go. & Label definitions need operational clarity. \\
\bottomrule
\end{tabularx}
\rowcolors{3}{}{}
\end{table}
